\documentclass[aspectratio=169,10pt]{beamer}

\usepackage{fontspec}
\usepackage{amsmath}
\usepackage{booktabs}
\usepackage{tikz}
\usepackage{xcolor}
\usetikzlibrary{positioning}

\setsansfont{DejaVu Sans}
\setmonofont{DejaVu Sans Mono}[Scale=MatchLowercase]

\definecolor{teal}{HTML}{087F82}
\definecolor{tealDark}{HTML}{05565B}
\definecolor{tealLight}{HTML}{DDF3F0}
\definecolor{mint}{HTML}{B7E5DE}
\definecolor{ink}{HTML}{12343B}
\definecolor{coral}{HTML}{FFB38A}

\usetheme{default}
\setbeamertemplate{navigation symbols}{}
\setbeamertemplate{blocks}[rounded][shadow=false]
\setbeamercolor{normal text}{fg=white,bg=teal}
\setbeamercolor{structure}{fg=mint}
\setbeamercolor{footline}{bg=tealDark,fg=white}
\setbeamercolor{frametitle}{fg=ink,bg=mint}
\setbeamercolor{block title}{fg=ink,bg=mint}
\setbeamercolor{block body}{fg=ink,bg=tealLight}
\setbeamercolor{alerted text}{fg=coral}
\setbeamerfont{frametitle}{size=\Large,series=\bfseries}
\setbeamerfont{title}{size=\LARGE,series=\bfseries}
\setbeamerfont{subtitle}{size=\large}
\setbeamertemplate{itemize item}{\color{coral}\small$\blacktriangleright$}
\setbeamertemplate{itemize subitem}{\color{mint}\scriptsize$\blacktriangleright$}

% PLACEHOLDER: thay hai gia tri nay truoc khi trinh bay.
\newcommand{\AuthorName}{Đỗ Thành Trung}
\newcommand{\UniversityName}{IAI, UET - VNU}
\newcommand{\DeckTitle}{Đánh giá các mô hình nhúng \\ cho tiếng Việt}

\setbeamertemplate{footline}{%
	\leavevmode\hbox{%
		\begin{beamercolorbox}[wd=.78\paperwidth,ht=2.8ex,dp=1.2ex,leftskip=1.2em]{author in head/foot}%
			\scriptsize\UniversityName
		\end{beamercolorbox}%
		\begin{beamercolorbox}[wd=.22\paperwidth,ht=2.8ex,dp=1.2ex,rightskip=1.2em plus 1fil]{date in head/foot}%
			\scriptsize\hfill\insertframenumber~/~\inserttotalframenumber
		\end{beamercolorbox}}%
}

\title{\DeckTitle}
\author{\AuthorName}
\institute{\UniversityName}
\date{\today}

\AtBeginSection[]{
    \begin{frame}{Nội dung trình bày}
        \setbeamertemplate{section in toc}[sections numbered]
        \setbeamerfont{section in toc}{series=\bfseries}
        \setbeamerfont{section in toc shaded}{series=\normalfont}
        \tableofcontents[currentsection]
    \end{frame}
}

\begin{document}

\begin{frame}[plain]
    \begin{tikzpicture}[remember picture,overlay]
        \fill[tealDark] (current page.south west) rectangle (current page.north east);
        \fill[mint,opacity=.12] (current page.north east) circle (3.2cm);
        \fill[coral,opacity=.16] (current page.south west) circle (2.4cm);
    \end{tikzpicture}
    \vfill
    {\usebeamerfont{title}\color{white}\DeckTitle\par}
    \vspace{0.35cm}
    {\usebeamerfont{subtitle}\color{mint}\insertsubtitle\par}
    \vspace{1.1cm}
    \begin{tabular}{@{}ll@{}}
        Người trình bày: & \textbf{\AuthorName}     \\
        Đơn vị:          & \textbf{\UniversityName}
    \end{tabular}
    \vfill
    \color{mint}\small\insertdate
\end{frame}

\begin{frame}{Nội dung trình bày}
    \setbeamertemplate{section in toc}[sections numbered]
    \tableofcontents
\end{frame}

\section{Tổng quan}
\begin{frame}{Phạm vi báo cáo}
    \begin{block}{Tổng quan}
        Tổng hợp một số bộ dữ liệu có nhiễu hay được sử dụng, một số mô hình \\ tạo vector nhúng và nêu hướng thực nghiệm tiếp theo.
    \end{block}
    \begin{enumerate}
        \item Các phương pháp phổ biến để đánh giá mô hình nhúng.
        \item Một số mô hình thường được sử dụng.
        \item Một số công việc kĩ thuật đã thực hiện.
        \item Đề xuất hướng đi tiếp theo.
    \end{enumerate}
\end{frame}

% \begin{frame}{Pipeline dữ liệu và benchmark STS}
%     \begin{columns}[T,onlytextwidth]
%         \column{.45\textwidth}
%         \begin{block}{Luồng xử lý}
%             \begin{enumerate}
%                 \item Đọc file local và lọc split nếu cần.
%                 \item Chuẩn hóa hai cột câu và một cột điểm vàng.
%                 \item Encode hai phía bằng cùng một service.
%                 \item Tính cosine similarity.
%                 \item Đo tương quan Spearman với đánh giá con người.
%             \end{enumerate}
%         \end{block}
%         \column{.51\textwidth}
%         \begin{block}{Cấu hình hiện tại}
%             \small
%             \texttt{dataset}: path, format, delimiter, split\\[0.15cm]
%             \texttt{benchmark}: text\_a, text\_b, score\\[0.15cm]
%             \texttt{model}: model\_name, adapter, device, batch\_size
%         \end{block}
%         \begin{alertblock}{Ý nghĩa}
%             Thay dataset hoặc model không làm thay đổi logic benchmark.
%         \end{alertblock}
%     \end{columns}
% \end{frame}

% \begin{frame}{Một số lưu ý trước khi chạy}
%     \begin{block}{Tách từ tiếng Việt}
%         Một số mô hình yêu cầu tách từ và thêm dấu ``\_'' bằng \texttt{pyvi} hoặc \texttt{underthesea} trước khi đưa vào embedding.
%     \end{block}
%     \begin{block}{Mô hình đa ngữ}
%         Các mô hình đa ngữ thường nhận câu nguyên bản, không cần tách từ thủ công.
%     \end{block}
%     \begin{alertblock}{Asymmetric Retrieval}
%         Với E5, BGE và các mô hình hiện đại tương tự, khi đánh giá Retrieval cần dùng đúng tiền tố: \texttt{query: } cho câu hỏi và \texttt{passage: } cho đoạn văn.
%     \end{alertblock}
% \end{frame}

\section{Các mô hình nhúng phổ biến}
\begin{frame}{Các mô hình tiếng Việt}
    \scriptsize
    \renewcommand{\arraystretch}{1.55}
    \begin{tabular}{@{}p{0.45\textwidth}p{0.12\textwidth}p{0.12\textwidth}p{0.20\textwidth}@{}}
        \textbf{Mô hình}                                      & \textbf{Tham số}    & \textbf{Số chiều ẩn} & \textbf{Dung lượng mô hình} \\
        \midrule
        \texttt{AITeamVN/Vietnamese\_Embedding}               & \textasciitilde568M & 1024                 & 2{,}17 GB                   \\
        \texttt{dangvantuan/vietnamese-embedding}             & \textasciitilde135M & 768                  & 515 MB                      \\
        \texttt{bkai-foundation-models/vietnamese-bi-encoder} & \textasciitilde135M & 768                  & 515 MB                      \\
        \texttt{contextboxai/halong\_embedding}               & \textasciitilde278M & 768                  & 1{,}06 GB                   \\
        \texttt{vinai/phobert-base-v2}                        & \textasciitilde135M & 768                  & 515 MB                      \\
        \bottomrule
    \end{tabular}
\end{frame}

\begin{frame}{Các mô hình đa ngữ}
    \scriptsize
    \renewcommand{\arraystretch}{1.55}
    \begin{tabular}{@{}p{0.45\textwidth}p{0.12\textwidth}p{0.12\textwidth}p{0.20\textwidth}@{}}
        \textbf{Mô hình}                           & \textbf{Tham số}    & \textbf{Số chiều ẩn} & \textbf{Dung lượng mô hình} \\
        \midrule
        \texttt{Alibaba-NLP/gte-multilingual-base} & \textasciitilde305M & 768                  & 582 MB                      \\
        \texttt{BAAI/bge-m3}                       & \textasciitilde568M & 1024                 & 2{,}17 GB                   \\
        \texttt{intfloat/multilingual-e5-base}     & \textasciitilde278M & 768                  & 1{,}06 GB                   \\
        \shortstack[l]{\texttt{sentence-transformers/}                                                                        \\\texttt{paraphrase-multilingual-mpnet-base-v2}} & \textasciitilde278M & 768 & 1{,}06 GB\\
        \bottomrule
    \end{tabular}
\end{frame}

\section{Một số thang đo dữ liệu có nhiễu}
\begin{frame}{Dữ liệu viết tắt, ký tự đặc biệt, sai chính tả, không dấu}
    \begin{block}{Dữ liệu mạng xã hội}
        \begin{itemize}
            \item \textbf{UIT-ViHSD}: bình luận Facebook, 3 mức độ thù địch, nhiều teencode.
            \item \textbf{UIT-VSMEC}: cảm xúc mạng xã hội, câu thiếu ngữ cảnh và viết tắt.
            \item \textbf{ViSpamReviews}: phát hiện đánh giá spam trên sàn thương mại điện tử.
            \item \textbf{VMDS}: cặp câu sai chính tả và câu đã sửa, đo độ bền vững.
        \end{itemize}
    \end{block}
    \begin{block}{Nhận xét}
        Dữ liệu lệch khỏi văn phong tiền huấn luyện, có thể dùng làm phép đo \\ khả năng khái quát.
    \end{block}
\end{frame}

\begin{frame}{Dữ liệu gồm nhiều thuật ngữ chuyên ngành khó}
    \begin{block}{Dữ liệu khoa học}
        \begin{itemize}
            \item \textbf{CQADupstack-VN (GreenNote)}: Truy xuất câu trả lời trên diễn đàn; câu hỏi và văn bản trộn chữ, code, log, công thức toán học, ...
            \item \textbf{Touche2020-VN và ArguAna-VN}: Truy xuất văn bản phản biện hoặc ủng hộ luận điểm. Văn bản có thể chỉ khác luận điểm đúng một chữ.
            \item \textbf{SciDocsRR-VN và SciFact-VN}: Các mẫu dương và âm gần giống hệt nhau, các văn bản gồm nhiều thuật ngữ chuyên ngành.
        \end{itemize}
    \end{block}
    \begin{block}{Nhận xét}
        Dữ liệu gồm nhiều từ ngữ hoặc văn bản chuyên ngành, mô hình nhúng chỉ tạo một vector khó biểu diễn ngữ cảnh.
    \end{block}
\end{frame}

\begin{frame}{Dữ liệu có cấu trúc văn bản khó}
    \begin{block}{Dữ liệu phản biện}
        \begin{itemize}
            \item \textbf{ViNLI}: Cặp tiền đề và giả thuyết được gán nhãn \textit{entailment}, \textit{neutral} hoặc \textit{contradiction}.
        \end{itemize}
    \end{block}
    \begin{block}{Nhận xét}
        Các văn bản đối lập vẫn có thể trùng lặp nhiều từ. Mô hình cần hiểu quan hệ và ngữ cảnh, không chỉ đếm mức độ giống bề mặt.
    \end{block}
\end{frame}

\section{Các thang đo đánh giá}
\begin{frame}{Các nhóm thang đo chính}
    \begin{itemize}
        \item \textbf{STS (Semantic Textual Similarity)}: Tính tương quan Pearson/Spearman.
        \item \textbf{Retrieval / Reranking}: Tính độ đo MRR@k, NDCG@k của kết quả truy hồi.
        \item \textbf{Classification}: Dùng mô hình phân loại đơn giản (Logistic Regression) trên không gian nhúng.
        \item \textbf{Pair Classification}: Đo mức độ đồng nghĩa, trung tính, trái nghĩa.
        \item \textbf{Clustering}: Sử dụng thước đo V-measure.
        \item \textbf{Intrinsic geometry}: Kiểm tra tính đẳng hướng trong không gian, tránh hiện tượng sụp đổ biểu diễn.

    \end{itemize}
\end{frame}

\begin{frame}{STS, Retrieval và các tác vụ giám sát}
    \begin{block}{Semantic Textual Similarity (STS)}
        So sánh và chấm điểm cho từng cặp câu, điểm số được đối chiếu với điểm con người bằng tương quan Pearson hoặc Spearman.
    \end{block}
    \begin{block}{Retrieval / Reranking}
        Tập trung so sánh vector nhúng của câu hỏi với tài liệu, thường dùng để đánh giá độ chính xác với ngữ cảnh dài.
    \end{block}
    \begin{block}{Classification}
        Trực giác: Nếu không gian nhúng tách biệt các chủ đề tốt, mô hình tuyến tính đơn giản cũng có thể đạt hiệu suất cao.
    \end{block}
\end{frame}

\begin{frame}{Clustering}
    \begin{block}{Clustering}
        Phân cụm với số cụm bằng số nhãn đúng, rồi đo chỉ số V-measure:
        \[
            V = 2\cdot\frac{h\cdot c}{h+c}
        \]
        \small
        \textbf{Homogeneity} $h$: Cao khi các điểm trong cùng cụm có cùng lớp.\\
        \textbf{Completeness} $c$: Cao khi các điểm trong cùng lớp được đặt trong cùng cụm.
    \end{block}
\end{frame}

\begin{frame}{Đánh giá nội tại không gian nhúng}
    \begin{block}{Tính liên kết: các câu cùng nghĩa nằm gần nhau}
        \[
            \mathcal{L}_{\mathrm{align}} = \mathbb{E}_{(x,x^+)}\left[\left\|f(x)-f(x^+)\right\|_2^2\right]
        \]
        Giá trị càng nhỏ càng tốt.
    \end{block}
    \begin{block}{Tính phân bố đều: vector phân bố đều trên mặt cầu đơn vị}
        \[
            \mathcal{L}_{\mathrm{uniform}} = \log\mathbb{E}_{x,y\sim p_{\mathrm{data}}}\left[e^{-2\left\|f(x)-f(y)\right\|_2^2}\right]
        \]
        Giá trị càng nhỏ/càng âm càng tốt, tránh sụp đổ biểu diễn hoặc \\ các vector chụm thành "hình nón hẹp".
    \end{block}
\end{frame}

\section{Các kiến trúc đã triển khai}
\begin{frame}{Các kiến trúc đã triển khai}
    \centering
    \begin{columns}[T,onlytextwidth]
        \column{.48\textwidth}
        \begin{block}{Thiết kế giao diện gọi mô hình}
            \begin{itemize}
                \item Viết giao diện để gọi mô hình, \\ sử dụng API của thư viện \texttt{sentence-transformers}.
                \item Hỗ trợ tiền xử lý cho từng mô hình trước khi chạy.
            \end{itemize}
        \end{block}
        \column{.48\textwidth}
        \begin{block}{Giao diện sử dụng dữ liệu}
            \begin{itemize}
                \item Hỗ trợ sử dụng dữ liệu từ API của \texttt{datasets}, gọi từ máy cá nhân hoặc URI.
                \item Chạy một số benchmark đơn giản trên các bộ dữ liệu nhỏ.
            \end{itemize}
        \end{block}
    \end{columns}
\end{frame}

\section{Lộ trình tiếp theo}
\begin{frame}{Tiến hành đánh giá kết quả các mô hình và \\ tìm các trường hợp biên}
    \begin{enumerate}
        \item Chuẩn hóa giao diện gọi các mô hình và bộ dữ liệu khác nhau.
        \item Chạy kết quả trên một số thang đo đánh giá khác nhau và \\ tổng hợp kết quả.
        \item Phân tích tính đẳng hướng, kiểm tra hiện tượng sụp đổ biểu diễn.
        \item Tìm và phân tích các trường hợp biên mà mô hình gặp khó.
    \end{enumerate}
\end{frame}

% \begin{frame}[plain]
%     \centering
%     \vfill
%     {\usebeamerfont{title}\color{white}Cảm ơn mọi người đã lắng nghe}\par
%     \vspace{0.4cm}
%     {\large\color{mint}\DeckTitle}\par
%     \vspace{0.8cm}
%     \textbf{\AuthorName}\\
%     \UniversityName
%     \vfill
% \end{frame}

\end{document}
